\documentclass[10pt]{article}
\usepackage[a4paper,landscape,margin=1.0cm]{geometry}
\usepackage{graphicx}
\usepackage{xcolor}
\usepackage{enumitem}
\usepackage{tcolorbox}
\definecolor{ofbTeal}{HTML}{0E6E6E}
\definecolor{ofbBlue}{HTML}{1F5C8B}
\definecolor{ofbInk}{HTML}{1A1A1A}
\definecolor{ofbMuted}{HTML}{5B5B5B}
\setlength{\parindent}{0pt}
\pagestyle{empty}
\begin{document}

\noindent\begin{minipage}[t]{0.585\textwidth}\vspace{0pt}
{\setlength{\fboxsep}{0pt}\setlength{\fboxrule}{0.4pt}\fbox{\includegraphics[width=\linewidth]{build/guion/g-01.png}}}
\end{minipage}\hfill
\begin{minipage}[t]{0.395\textwidth}\vspace{2pt}
{\footnotesize\color{ofbMuted}Slide 1 / 60 \quad $\sim$30s \quad cum 0:30}\hfill{\footnotesize\bfseries\color{ofbTeal}CORE}\\[2pt]
{\large\bfseries\color{ofbTeal} Title}\\[6pt]
{\small\itshape\color{ofbMuted}Say:}\\[2pt]
\begin{itemize}[leftmargin=1.1em,itemsep=3pt,topsep=2pt]
\item Good morning, everyone. I am Jorge Mayorga, from Universidad de los Andes.
\item I want to start with a deceptively simple question: in a grid run by inverters, which frequency estimator should we actually trust?
\item The honest answer is, it depends. This talk makes 'it depends' precise.
\end{itemize}
\vspace{2pt}{\small\color{ofbBlue}$\rightarrow$ \itshape Let me start with why this is a problem worth solving.}\\[2pt]
\vspace{3pt}
\begin{tcolorbox}[colback=ofbTeal!8,colframe=ofbTeal,boxrule=0.6pt,arc=2pt,left=4pt,right=4pt,top=3pt,bottom=3pt]
{\small\bfseries Key line: }{\small There is no universal best estimator; the right choice depends on what you measure.}
\end{tcolorbox}
\end{minipage}
\clearpage
\noindent\begin{minipage}[t]{0.585\textwidth}\vspace{0pt}
{\setlength{\fboxsep}{0pt}\setlength{\fboxrule}{0.4pt}\fbox{\includegraphics[width=\linewidth]{build/guion/g-02.png}}}
\end{minipage}\hfill
\begin{minipage}[t]{0.395\textwidth}\vspace{2pt}
{\footnotesize\color{ofbMuted}Slide 2 / 60 \quad $\sim$6s \quad cum 0:36}\hfill{\footnotesize\bfseries\color{ofbMuted}QUICK}\\[2pt]
{\large\bfseries\color{ofbMuted} Motivation (section)}\\[6pt]
\vspace{3pt}
\begin{tcolorbox}[colback=ofbMuted!8,colframe=ofbMuted,boxrule=0.6pt,arc=2pt,left=4pt,right=4pt,top=3pt,bottom=3pt]
{\small\bfseries Key line: }{\small Why this matters.}
\end{tcolorbox}
\end{minipage}
\clearpage
\noindent\begin{minipage}[t]{0.585\textwidth}\vspace{0pt}
{\setlength{\fboxsep}{0pt}\setlength{\fboxrule}{0.4pt}\fbox{\includegraphics[width=\linewidth]{build/guion/g-03.png}}}
\end{minipage}\hfill
\begin{minipage}[t]{0.395\textwidth}\vspace{2pt}
{\footnotesize\color{ofbMuted}Slide 3 / 60 \quad $\sim$35s \quad cum 1:11}\hfill{\footnotesize\bfseries\color{ofbTeal}CORE}\\[2pt]
{\large\bfseries\color{ofbTeal} Field events define the stress model}\\[6pt]
{\small\itshape\color{ofbMuted}Say:}\\[2pt]
\begin{itemize}[leftmargin=1.1em,itemsep=3pt,topsep=2pt]
\item Real inverter disturbances are nothing like a clean sine wave.
\item In events like Blue Cut, Odessa, and the Kauai oscillation, phase jumps, fast RoCoF, harmonics and noise all arrive together.
\item So we let these real events define the stress cases our benchmark has to cover.
\end{itemize}
\vspace{2pt}{\small\color{ofbBlue}$\rightarrow$ \itshape And the tool we use to see these events also shapes what we conclude.}\\[2pt]
\vspace{3pt}
\begin{tcolorbox}[colback=ofbTeal!8,colframe=ofbTeal,boxrule=0.6pt,arc=2pt,left=4pt,right=4pt,top=3pt,bottom=3pt]
{\small\bfseries Key line: }{\small Field events, not textbook signals, set the bar.}
\end{tcolorbox}
\end{minipage}
\clearpage
\noindent\begin{minipage}[t]{0.585\textwidth}\vspace{0pt}
{\setlength{\fboxsep}{0pt}\setlength{\fboxrule}{0.4pt}\fbox{\includegraphics[width=\linewidth]{build/guion/g-04.png}}}
\end{minipage}\hfill
\begin{minipage}[t]{0.395\textwidth}\vspace{2pt}
{\footnotesize\color{ofbMuted}Slide 4 / 60 \quad $\sim$30s \quad cum 1:41}\hfill{\footnotesize\bfseries\color{ofbTeal}CORE}\\[2pt]
{\large\bfseries\color{ofbTeal} PMU streams show consequences}\\[6pt]
{\small\itshape\color{ofbMuted}Say:}\\[2pt]
\begin{itemize}[leftmargin=1.1em,itemsep=3pt,topsep=2pt]
\item A standard PMU gives us frequency and RoCoF.
\item But during an inverter event, those numbers show the consequence, not always the cause, and the estimator inside the PMU decides what we see.
\end{itemize}
\vspace{2pt}{\small\color{ofbBlue}$\rightarrow$ \itshape So before we can analyze events, we have to ask which estimator to trust.}\\[2pt]
\vspace{3pt}
\begin{tcolorbox}[colback=ofbTeal!8,colframe=ofbTeal,boxrule=0.6pt,arc=2pt,left=4pt,right=4pt,top=3pt,bottom=3pt]
{\small\bfseries Key line: }{\small Choosing the estimator is itself a measurement decision.}
\end{tcolorbox}
\end{minipage}
\clearpage
\noindent\begin{minipage}[t]{0.585\textwidth}\vspace{0pt}
{\setlength{\fboxsep}{0pt}\setlength{\fboxrule}{0.4pt}\fbox{\includegraphics[width=\linewidth]{build/guion/g-05.png}}}
\end{minipage}\hfill
\begin{minipage}[t]{0.395\textwidth}\vspace{2pt}
{\footnotesize\color{ofbMuted}Slide 5 / 60 \quad $\sim$28s \quad cum 2:09}\hfill{\footnotesize\bfseries\color{ofbTeal}CORE}\\[2pt]
{\large\bfseries\color{ofbTeal} Research problem}\\[6pt]
{\small\itshape\color{ofbMuted}Say:}\\[2pt]
\begin{itemize}[leftmargin=1.1em,itemsep=3pt,topsep=2pt]
\item And here is the gap: there is no agreed way to compare these estimators under inverter-specific stress.
\item Compliance tests use clean signals, so they never tell us what happens on a real composite event.
\end{itemize}
\vspace{2pt}{\small\color{ofbBlue}$\rightarrow$ \itshape Our answer to that gap is one clear claim.}\\[2pt]
\vspace{3pt}
\begin{tcolorbox}[colback=ofbTeal!8,colframe=ofbTeal,boxrule=0.6pt,arc=2pt,left=4pt,right=4pt,top=3pt,bottom=3pt]
{\small\bfseries Key line: }{\small We need a fair, inverter-relevant way to compare estimators.}
\end{tcolorbox}
\end{minipage}
\clearpage
\noindent\begin{minipage}[t]{0.585\textwidth}\vspace{0pt}
{\setlength{\fboxsep}{0pt}\setlength{\fboxrule}{0.4pt}\fbox{\includegraphics[width=\linewidth]{build/guion/g-06.png}}}
\end{minipage}\hfill
\begin{minipage}[t]{0.395\textwidth}\vspace{2pt}
{\footnotesize\color{ofbMuted}Slide 6 / 60 \quad $\sim$25s \quad cum 2:34}\hfill{\footnotesize\bfseries\color{ofbTeal}CORE}\\[2pt]
{\large\bfseries\color{ofbTeal} Operating claim}\\[6pt]
{\small\itshape\color{ofbMuted}Say:}\\[2pt]
\begin{itemize}[leftmargin=1.1em,itemsep=3pt,topsep=2pt]
\item Our claim is that estimator selection has to be conditional.
\item It depends on the disturbance type, the metric you care about, and your compute budget.
\end{itemize}
\vspace{2pt}{\small\color{ofbBlue}$\rightarrow$ \itshape To support that claim we built three things.}\\[2pt]
\vspace{3pt}
\begin{tcolorbox}[colback=ofbTeal!8,colframe=ofbTeal,boxrule=0.6pt,arc=2pt,left=4pt,right=4pt,top=3pt,bottom=3pt]
{\small\bfseries Key line: }{\small There is no single best estimator, only a best estimator for a given job.}
\end{tcolorbox}
\end{minipage}
\clearpage
\noindent\begin{minipage}[t]{0.585\textwidth}\vspace{0pt}
{\setlength{\fboxsep}{0pt}\setlength{\fboxrule}{0.4pt}\fbox{\includegraphics[width=\linewidth]{build/guion/g-07.png}}}
\end{minipage}\hfill
\begin{minipage}[t]{0.395\textwidth}\vspace{2pt}
{\footnotesize\color{ofbMuted}Slide 7 / 60 \quad $\sim$30s \quad cum 3:04}\hfill{\footnotesize\bfseries\color{ofbTeal}CORE}\\[2pt]
{\large\bfseries\color{ofbTeal} Technical contributions}\\[6pt]
{\small\itshape\color{ofbMuted}Say:}\\[2pt]
\begin{itemize}[leftmargin=1.1em,itemsep=3pt,topsep=2pt]
\item First, an open benchmark platform with known-truth scenarios.
\item Second, a common execution contract so all eighteen estimators are tested identically.
\item And third, a metric vector that keeps accuracy, trip-risk, latency and cost separate.
\end{itemize}
\vspace{2pt}{\small\color{ofbBlue}$\rightarrow$ \itshape Let me give you a sense of the scale.}\\[2pt]
\vspace{3pt}
\begin{tcolorbox}[colback=ofbTeal!8,colframe=ofbTeal,boxrule=0.6pt,arc=2pt,left=4pt,right=4pt,top=3pt,bottom=3pt]
{\small\bfseries Key line: }{\small Open platform, common contract, and multi-metric evidence.}
\end{tcolorbox}
\end{minipage}
\clearpage
\noindent\begin{minipage}[t]{0.585\textwidth}\vspace{0pt}
{\setlength{\fboxsep}{0pt}\setlength{\fboxrule}{0.4pt}\fbox{\includegraphics[width=\linewidth]{build/guion/g-08.png}}}
\end{minipage}\hfill
\begin{minipage}[t]{0.395\textwidth}\vspace{2pt}
{\footnotesize\color{ofbMuted}Slide 8 / 60 \quad $\sim$25s \quad cum 3:29}\hfill{\footnotesize\bfseries\color{ofbTeal}CORE}\\[2pt]
{\large\bfseries\color{ofbTeal} Experiment scale}\\[6pt]
{\small\itshape\color{ofbMuted}Say:}\\[2pt]
\begin{itemize}[leftmargin=1.1em,itemsep=3pt,topsep=2pt]
\item Eighteen estimators, thirty-three scenarios, five families of stress, with Monte-Carlo repetitions.
\item On top of that, a severity-sweep engine we call ATLAS, and everything is reproducible from artifacts.
\end{itemize}
\vspace{2pt}{\small\color{ofbBlue}$\rightarrow$ \itshape And the logic for choosing between all these methods fits in one picture.}\\[2pt]
\vspace{3pt}
\begin{tcolorbox}[colback=ofbTeal!8,colframe=ofbTeal,boxrule=0.6pt,arc=2pt,left=4pt,right=4pt,top=3pt,bottom=3pt]
{\small\bfseries Key line: }{\small Large, factorial, and fully reproducible.}
\end{tcolorbox}
\end{minipage}
\clearpage
\noindent\begin{minipage}[t]{0.585\textwidth}\vspace{0pt}
{\setlength{\fboxsep}{0pt}\setlength{\fboxrule}{0.4pt}\fbox{\includegraphics[width=\linewidth]{build/guion/g-09.png}}}
\end{minipage}\hfill
\begin{minipage}[t]{0.395\textwidth}\vspace{2pt}
{\footnotesize\color{ofbMuted}Slide 9 / 60 \quad $\sim$30s \quad cum 3:59}\hfill{\footnotesize\bfseries\color{ofbBlue}FIG}\\[2pt]
{\large\bfseries\color{ofbBlue} Estimator selection logic (build-up)}\\[6pt]
{\small\itshape\color{ofbMuted}Say:}\\[2pt]
\begin{itemize}[leftmargin=1.1em,itemsep=3pt,topsep=2pt]
\item Selection comes down to three questions: which event, which metric, and which compute budget.
\item Watch as each one is revealed, because each can point to a different method.
\end{itemize}
\vspace{2pt}{\small\color{ofbBlue}$\rightarrow$ \itshape With that logic in mind, let me show you the platform itself.}\\[2pt]
\vspace{3pt}
\begin{tcolorbox}[colback=ofbBlue!8,colframe=ofbBlue,boxrule=0.6pt,arc=2pt,left=4pt,right=4pt,top=3pt,bottom=3pt]
{\small\bfseries Key line: }{\small This one diagram is the whole talk in miniature.}
\end{tcolorbox}
\end{minipage}
\clearpage
\noindent\begin{minipage}[t]{0.585\textwidth}\vspace{0pt}
{\setlength{\fboxsep}{0pt}\setlength{\fboxrule}{0.4pt}\fbox{\includegraphics[width=\linewidth]{build/guion/g-10.png}}}
\end{minipage}\hfill
\begin{minipage}[t]{0.395\textwidth}\vspace{2pt}
{\footnotesize\color{ofbMuted}Slide 10 / 60 \quad $\sim$6s \quad cum 4:05}\hfill{\footnotesize\bfseries\color{ofbMuted}QUICK}\\[2pt]
{\large\bfseries\color{ofbMuted} Benchmark and platform (section)}\\[6pt]
\vspace{3pt}
\begin{tcolorbox}[colback=ofbMuted!8,colframe=ofbMuted,boxrule=0.6pt,arc=2pt,left=4pt,right=4pt,top=3pt,bottom=3pt]
{\small\bfseries Key line: }{\small First, the benchmark and the platform.}
\end{tcolorbox}
\end{minipage}
\clearpage
\noindent\begin{minipage}[t]{0.585\textwidth}\vspace{0pt}
{\setlength{\fboxsep}{0pt}\setlength{\fboxrule}{0.4pt}\fbox{\includegraphics[width=\linewidth]{build/guion/g-11.png}}}
\end{minipage}\hfill
\begin{minipage}[t]{0.395\textwidth}\vspace{2pt}
{\footnotesize\color{ofbMuted}Slide 11 / 60 \quad $\sim$25s \quad cum 4:30}\hfill{\footnotesize\bfseries\color{ofbTeal}CORE}\\[2pt]
{\large\bfseries\color{ofbTeal} Benchmark literature gap}\\[6pt]
{\small\itshape\color{ofbMuted}Say:}\\[2pt]
\begin{itemize}[leftmargin=1.1em,itemsep=3pt,topsep=2pt]
\item Prior work usually compares a handful of estimators on one or two cases.
\item Nobody has covered the full estimator space against composite inverter stress with reproducible metrics.
\end{itemize}
\vspace{2pt}{\small\color{ofbBlue}$\rightarrow$ \itshape So we built exactly that.}\\[2pt]
\vspace{3pt}
\begin{tcolorbox}[colback=ofbTeal!8,colframe=ofbTeal,boxrule=0.6pt,arc=2pt,left=4pt,right=4pt,top=3pt,bottom=3pt]
{\small\bfseries Key line: }{\small That blank space on the map is what we fill.}
\end{tcolorbox}
\end{minipage}
\clearpage
\noindent\begin{minipage}[t]{0.585\textwidth}\vspace{0pt}
{\setlength{\fboxsep}{0pt}\setlength{\fboxrule}{0.4pt}\fbox{\includegraphics[width=\linewidth]{build/guion/g-12.png}}}
\end{minipage}\hfill
\begin{minipage}[t]{0.395\textwidth}\vspace{2pt}
{\footnotesize\color{ofbMuted}Slide 12 / 60 \quad $\sim$25s \quad cum 4:55}\hfill{\footnotesize\bfseries\color{ofbTeal}CORE}\\[2pt]
{\large\bfseries\color{ofbTeal} Benchmark comparison axes}\\[6pt]
{\small\itshape\color{ofbMuted}Say:}\\[2pt]
\begin{itemize}[leftmargin=1.1em,itemsep=3pt,topsep=2pt]
\item We compare on four axes at once: accuracy, robustness, computational cost, and protection trip exposure.
\item The key point is that a method can win one axis and lose another.
\end{itemize}
\vspace{2pt}{\small\color{ofbBlue}$\rightarrow$ \itshape Here is the platform that measures all four.}\\[2pt]
\vspace{3pt}
\begin{tcolorbox}[colback=ofbTeal!8,colframe=ofbTeal,boxrule=0.6pt,arc=2pt,left=4pt,right=4pt,top=3pt,bottom=3pt]
{\small\bfseries Key line: }{\small One number hides the trade-off; four axes reveal it.}
\end{tcolorbox}
\end{minipage}
\clearpage
\noindent\begin{minipage}[t]{0.585\textwidth}\vspace{0pt}
{\setlength{\fboxsep}{0pt}\setlength{\fboxrule}{0.4pt}\fbox{\includegraphics[width=\linewidth]{build/guion/g-13.png}}}
\end{minipage}\hfill
\begin{minipage}[t]{0.395\textwidth}\vspace{2pt}
{\footnotesize\color{ofbMuted}Slide 13 / 60 \quad $\sim$28s \quad cum 5:23}\hfill{\footnotesize\bfseries\color{ofbBlue}FIG}\\[2pt]
{\large\bfseries\color{ofbBlue} OpenFreqBench platform}\\[6pt]
{\small\itshape\color{ofbMuted}Say:}\\[2pt]
\begin{itemize}[leftmargin=1.1em,itemsep=3pt,topsep=2pt]
\item OpenFreqBench has four parts: a scenario factory with known truth, a uniform estimator interface, a locked metric engine, and traceable artifacts.
\end{itemize}
\vspace{2pt}{\small\color{ofbBlue}$\rightarrow$ \itshape Let me show you how we keep the comparison fair.}\\[2pt]
\vspace{3pt}
\begin{tcolorbox}[colback=ofbBlue!8,colframe=ofbBlue,boxrule=0.6pt,arc=2pt,left=4pt,right=4pt,top=3pt,bottom=3pt]
{\small\bfseries Key line: }{\small Known truth in, locked metrics out, everything traced.}
\end{tcolorbox}
\end{minipage}
\clearpage
\noindent\begin{minipage}[t]{0.585\textwidth}\vspace{0pt}
{\setlength{\fboxsep}{0pt}\setlength{\fboxrule}{0.4pt}\fbox{\includegraphics[width=\linewidth]{build/guion/g-14.png}}}
\end{minipage}\hfill
\begin{minipage}[t]{0.395\textwidth}\vspace{2pt}
{\footnotesize\color{ofbMuted}Slide 14 / 60 \quad $\sim$22s \quad cum 5:45}\hfill{\footnotesize\bfseries\color{ofbTeal}CORE}\\[2pt]
{\large\bfseries\color{ofbTeal} Experimental method (1)}\\[6pt]
{\small\itshape\color{ofbMuted}Say:}\\[2pt]
\begin{itemize}[leftmargin=1.1em,itemsep=3pt,topsep=2pt]
\item Fairness is the whole point: every estimator sees the exact same input trace, the same random seeds, and the same warm-up handling.
\end{itemize}
\vspace{2pt}{\small\color{ofbBlue}$\rightarrow$ \itshape And we are equally strict about tuning and timing.}\\[2pt]
\vspace{3pt}
\begin{tcolorbox}[colback=ofbTeal!8,colframe=ofbTeal,boxrule=0.6pt,arc=2pt,left=4pt,right=4pt,top=3pt,bottom=3pt]
{\small\bfseries Key line: }{\small Fairness is enforced by construction, not by trust.}
\end{tcolorbox}
\end{minipage}
\clearpage
\noindent\begin{minipage}[t]{0.585\textwidth}\vspace{0pt}
{\setlength{\fboxsep}{0pt}\setlength{\fboxrule}{0.4pt}\fbox{\includegraphics[width=\linewidth]{build/guion/g-15.png}}}
\end{minipage}\hfill
\begin{minipage}[t]{0.395\textwidth}\vspace{2pt}
{\footnotesize\color{ofbMuted}Slide 15 / 60 \quad $\sim$22s \quad cum 6:07}\hfill{\footnotesize\bfseries\color{ofbBlue}FIG}\\[2pt]
{\large\bfseries\color{ofbBlue} Experimental method (2)}\\[6pt]
{\small\itshape\color{ofbMuted}Say:}\\[2pt]
\begin{itemize}[leftmargin=1.1em,itemsep=3pt,topsep=2pt]
\item We tune under a fixed policy, we drop the warm-up region, and we measure CPU as a real output.
\end{itemize}
\vspace{2pt}{\small\color{ofbBlue}$\rightarrow$ \itshape Putting it together, this is the full pipeline.}\\[2pt]
\vspace{3pt}
\begin{tcolorbox}[colback=ofbBlue!8,colframe=ofbBlue,boxrule=0.6pt,arc=2pt,left=4pt,right=4pt,top=3pt,bottom=3pt]
{\small\bfseries Key line: }{\small Nothing is hand-picked to flatter one method.}
\end{tcolorbox}
\end{minipage}
\clearpage
\noindent\begin{minipage}[t]{0.585\textwidth}\vspace{0pt}
{\setlength{\fboxsep}{0pt}\setlength{\fboxrule}{0.4pt}\fbox{\includegraphics[width=\linewidth]{build/guion/g-16.png}}}
\end{minipage}\hfill
\begin{minipage}[t]{0.395\textwidth}\vspace{2pt}
{\footnotesize\color{ofbMuted}Slide 16 / 60 \quad $\sim$22s \quad cum 6:29}\hfill{\footnotesize\bfseries\color{ofbBlue}FIG}\\[2pt]
{\large\bfseries\color{ofbBlue} Full benchmark workflow}\\[6pt]
{\small\itshape\color{ofbMuted}Say:}\\[2pt]
\begin{itemize}[leftmargin=1.1em,itemsep=3pt,topsep=2pt]
\item Scenario, Monte-Carlo, estimator, locked metrics, then artifacts.
\item Every single figure in this talk traces back to those artifacts.
\end{itemize}
\vspace{2pt}{\small\color{ofbBlue}$\rightarrow$ \itshape Let me show you the scenarios themselves.}\\[2pt]
\vspace{3pt}
\begin{tcolorbox}[colback=ofbBlue!8,colframe=ofbBlue,boxrule=0.6pt,arc=2pt,left=4pt,right=4pt,top=3pt,bottom=3pt]
{\small\bfseries Key line: }{\small One reproducible pipeline, end to end.}
\end{tcolorbox}
\end{minipage}
\clearpage
\noindent\begin{minipage}[t]{0.585\textwidth}\vspace{0pt}
{\setlength{\fboxsep}{0pt}\setlength{\fboxrule}{0.4pt}\fbox{\includegraphics[width=\linewidth]{build/guion/g-17.png}}}
\end{minipage}\hfill
\begin{minipage}[t]{0.395\textwidth}\vspace{2pt}
{\footnotesize\color{ofbMuted}Slide 17 / 60 \quad $\sim$25s \quad cum 6:54}\hfill{\footnotesize\bfseries\color{ofbTeal}CORE}\\[2pt]
{\large\bfseries\color{ofbTeal} Standard stresses: controlled waveform}\\[6pt]
{\small\itshape\color{ofbMuted}Say:}\\[2pt]
\begin{itemize}[leftmargin=1.1em,itemsep=3pt,topsep=2pt]
\item We start with the standard stresses: a step, a ramp, and modulation.
\item Because the waveform is fully controlled, we know the true frequency at every single sample.
\end{itemize}
\vspace{2pt}{\small\color{ofbBlue}$\rightarrow$ \itshape Here is what those look like.}\\[2pt]
\vspace{3pt}
\begin{tcolorbox}[colback=ofbTeal!8,colframe=ofbTeal,boxrule=0.6pt,arc=2pt,left=4pt,right=4pt,top=3pt,bottom=3pt]
{\small\bfseries Key line: }{\small Known truth is what lets us score the error exactly.}
\end{tcolorbox}
\end{minipage}
\clearpage
\noindent\begin{minipage}[t]{0.585\textwidth}\vspace{0pt}
{\setlength{\fboxsep}{0pt}\setlength{\fboxrule}{0.4pt}\fbox{\includegraphics[width=\linewidth]{build/guion/g-18.png}}}
\end{minipage}\hfill
\begin{minipage}[t]{0.395\textwidth}\vspace{2pt}
{\footnotesize\color{ofbMuted}Slide 18 / 60 \quad $\sim$30s \quad cum 7:24}\hfill{\footnotesize\bfseries\color{ofbBlue}FIG}\\[2pt]
{\large\bfseries\color{ofbBlue} Scenario suite I (dashboard)}\\[6pt]
{\small\itshape\color{ofbMuted}Say:}\\[2pt]
\begin{itemize}[leftmargin=1.1em,itemsep=3pt,topsep=2pt]
\item Voltage on top, true frequency on the bottom, one column per scenario.
\item Notice that for the magnitude step and the AM case, the frequency stays flat at sixty hertz, yet the estimators still react.
\end{itemize}
\vspace{2pt}{\small\color{ofbBlue}$\rightarrow$ \itshape Now we make it realistic.}\\[2pt]
\vspace{3pt}
\begin{tcolorbox}[colback=ofbBlue!8,colframe=ofbBlue,boxrule=0.6pt,arc=2pt,left=4pt,right=4pt,top=3pt,bottom=3pt]
{\small\bfseries Key line: }{\small The stress is not always a frequency event, but estimators react anyway.}
\end{tcolorbox}
\end{minipage}
\clearpage
\noindent\begin{minipage}[t]{0.585\textwidth}\vspace{0pt}
{\setlength{\fboxsep}{0pt}\setlength{\fboxrule}{0.4pt}\fbox{\includegraphics[width=\linewidth]{build/guion/g-19.png}}}
\end{minipage}\hfill
\begin{minipage}[t]{0.395\textwidth}\vspace{2pt}
{\footnotesize\color{ofbMuted}Slide 19 / 60 \quad $\sim$25s \quad cum 7:49}\hfill{\footnotesize\bfseries\color{ofbTeal}CORE}\\[2pt]
{\large\bfseries\color{ofbTeal} Composite IBR stresses}\\[6pt]
{\small\itshape\color{ofbMuted}Say:}\\[2pt]
\begin{itemize}[leftmargin=1.1em,itemsep=3pt,topsep=2pt]
\item Now we stack everything: harmonics, phase jumps, RoCoF segments and ringdown, all in one waveform.
\item This is what a real inverter-dominated fault actually looks like.
\end{itemize}
\vspace{2pt}{\small\color{ofbBlue}$\rightarrow$ \itshape Here is that composite suite.}\\[2pt]
\vspace{3pt}
\begin{tcolorbox}[colback=ofbTeal!8,colframe=ofbTeal,boxrule=0.6pt,arc=2pt,left=4pt,right=4pt,top=3pt,bottom=3pt]
{\small\bfseries Key line: }{\small Composite stress is the realistic, and the revealing, case.}
\end{tcolorbox}
\end{minipage}
\clearpage
\noindent\begin{minipage}[t]{0.585\textwidth}\vspace{0pt}
{\setlength{\fboxsep}{0pt}\setlength{\fboxrule}{0.4pt}\fbox{\includegraphics[width=\linewidth]{build/guion/g-20.png}}}
\end{minipage}\hfill
\begin{minipage}[t]{0.395\textwidth}\vspace{2pt}
{\footnotesize\color{ofbMuted}Slide 20 / 60 \quad $\sim$25s \quad cum 8:14}\hfill{\footnotesize\bfseries\color{ofbBlue}FIG}\\[2pt]
{\large\bfseries\color{ofbBlue} Scenario suite II (dashboard)}\\[6pt]
{\small\itshape\color{ofbMuted}Say:}\\[2pt]
\begin{itemize}[leftmargin=1.1em,itemsep=3pt,topsep=2pt]
\item Same layout, voltage on top, frequency below, but these cases are deliberately harder than any compliance test.
\end{itemize}
\vspace{2pt}{\small\color{ofbBlue}$\rightarrow$ \itshape Before the results, a quick word on the methods and the metrics.}\\[2pt]
\vspace{3pt}
\begin{tcolorbox}[colback=ofbBlue!8,colframe=ofbBlue,boxrule=0.6pt,arc=2pt,left=4pt,right=4pt,top=3pt,bottom=3pt]
{\small\bfseries Key line: }{\small If a method survives here, it has earned our trust.}
\end{tcolorbox}
\end{minipage}
\clearpage
\noindent\begin{minipage}[t]{0.585\textwidth}\vspace{0pt}
{\setlength{\fboxsep}{0pt}\setlength{\fboxrule}{0.4pt}\fbox{\includegraphics[width=\linewidth]{build/guion/g-21.png}}}
\end{minipage}\hfill
\begin{minipage}[t]{0.395\textwidth}\vspace{2pt}
{\footnotesize\color{ofbMuted}Slide 21 / 60 \quad $\sim$28s \quad cum 8:42}\hfill{\footnotesize\bfseries\color{ofbTeal}CORE}\\[2pt]
{\large\bfseries\color{ofbTeal} Estimator families}\\[6pt]
{\small\itshape\color{ofbMuted}Say:}\\[2pt]
\begin{itemize}[leftmargin=1.1em,itemsep=3pt,topsep=2pt]
\item The eighteen estimators fall into families: loop-based PLLs, window and spectral methods, model-based Kalman filters, adaptive, and data-driven.
\item We analyze by family, because the family is what explains the behavior.
\end{itemize}
\vspace{2pt}{\small\color{ofbBlue}$\rightarrow$ \itshape And these are the numbers we score them on.}\\[2pt]
\vspace{3pt}
\begin{tcolorbox}[colback=ofbTeal!8,colframe=ofbTeal,boxrule=0.6pt,arc=2pt,left=4pt,right=4pt,top=3pt,bottom=3pt]
{\small\bfseries Key line: }{\small Eighteen estimators, five families, one fair contract.}
\end{tcolorbox}
\end{minipage}
\clearpage
\noindent\begin{minipage}[t]{0.585\textwidth}\vspace{0pt}
{\setlength{\fboxsep}{0pt}\setlength{\fboxrule}{0.4pt}\fbox{\includegraphics[width=\linewidth]{build/guion/g-22.png}}}
\end{minipage}\hfill
\begin{minipage}[t]{0.395\textwidth}\vspace{2pt}
{\footnotesize\color{ofbMuted}Slide 22 / 60 \quad $\sim$28s \quad cum 9:10}\hfill{\footnotesize\bfseries\color{ofbTeal}CORE}\\[2pt]
{\large\bfseries\color{ofbTeal} Metrics and decision variables}\\[6pt]
{\small\itshape\color{ofbMuted}Say:}\\[2pt]
\begin{itemize}[leftmargin=1.1em,itemsep=3pt,topsep=2pt]
\item We track RMSE for average error, peak error, and a trip-risk time, which counts how long the error stays above half a hertz.
\item Plus settling time and CPU cost.
\end{itemize}
\vspace{2pt}{\small\color{ofbBlue}$\rightarrow$ \itshape Now let me show you what the data says.}\\[2pt]
\vspace{3pt}
\begin{tcolorbox}[colback=ofbTeal!8,colframe=ofbTeal,boxrule=0.6pt,arc=2pt,left=4pt,right=4pt,top=3pt,bottom=3pt]
{\small\bfseries Key line: }{\small Half a hertz is our protection-oriented risk proxy.}
\end{tcolorbox}
\end{minipage}
\clearpage
\noindent\begin{minipage}[t]{0.585\textwidth}\vspace{0pt}
{\setlength{\fboxsep}{0pt}\setlength{\fboxrule}{0.4pt}\fbox{\includegraphics[width=\linewidth]{build/guion/g-23.png}}}
\end{minipage}\hfill
\begin{minipage}[t]{0.395\textwidth}\vspace{2pt}
{\footnotesize\color{ofbMuted}Slide 23 / 60 \quad $\sim$6s \quad cum 9:16}\hfill{\footnotesize\bfseries\color{ofbMuted}QUICK}\\[2pt]
{\large\bfseries\color{ofbMuted} Core evidence (section)}\\[6pt]
\vspace{3pt}
\begin{tcolorbox}[colback=ofbMuted!8,colframe=ofbMuted,boxrule=0.6pt,arc=2pt,left=4pt,right=4pt,top=3pt,bottom=3pt]
{\small\bfseries Key line: }{\small The core evidence: five scenarios, from standard to composite.}
\end{tcolorbox}
\end{minipage}
\clearpage
\noindent\begin{minipage}[t]{0.585\textwidth}\vspace{0pt}
{\setlength{\fboxsep}{0pt}\setlength{\fboxrule}{0.4pt}\fbox{\includegraphics[width=\linewidth]{build/guion/g-24.png}}}
\end{minipage}\hfill
\begin{minipage}[t]{0.395\textwidth}\vspace{2pt}
{\footnotesize\color{ofbMuted}Slide 24 / 60 \quad $\sim$40s \quad cum 9:56}\hfill{\footnotesize\bfseries\color{ofbTeal}CORE}\\[2pt]
{\large\bfseries\color{ofbTeal} Standard scenario results (build-up)}\\[6pt]
{\small\itshape\color{ofbMuted}Say:}\\[2pt]
\begin{itemize}[leftmargin=1.1em,itemsep=3pt,topsep=2pt]
\item On the standard cases, most methods are excellent, all below fifty millihertz.
\item CLICK. But watch the ramp: the textbook EKF and IpDFT intermittently blow up, to thirty-seven and fourteen hertz.
\end{itemize}
\vspace{2pt}{\small\color{ofbBlue}$\rightarrow$ \itshape So a phase jump should be even worse, right? Let me show you.}\\[2pt]
\vspace{3pt}
\begin{tcolorbox}[colback=ofbTeal!8,colframe=ofbTeal,boxrule=0.6pt,arc=2pt,left=4pt,right=4pt,top=3pt,bottom=3pt]
{\small\bfseries Key line: }{\small Robustness, not point accuracy, is what separates the families.}
\end{tcolorbox}
\end{minipage}
\clearpage
\noindent\begin{minipage}[t]{0.585\textwidth}\vspace{0pt}
{\setlength{\fboxsep}{0pt}\setlength{\fboxrule}{0.4pt}\fbox{\includegraphics[width=\linewidth]{build/guion/g-25.png}}}
\end{minipage}\hfill
\begin{minipage}[t]{0.395\textwidth}\vspace{2pt}
{\footnotesize\color{ofbMuted}Slide 25 / 60 \quad $\sim$35s \quad cum 10:31}\hfill{\footnotesize\bfseries\color{ofbTeal}CORE}\\[2pt]
{\large\bfseries\color{ofbTeal} Scenario D: phase jump separates families}\\[6pt]
{\small\itshape\color{ofbMuted}Say:}\\[2pt]
\begin{itemize}[leftmargin=1.1em,itemsep=3pt,topsep=2pt]
\item A clean sixty-degree phase jump cleanly splits the families.
\item The state-space filters and the PLLs recover with zero trip exposure.
\item The windowed and spectral methods carry over a second of exposure.
\end{itemize}
\vspace{2pt}{\small\color{ofbBlue}$\rightarrow$ \itshape But a single jump is not the real challenge.}\\[2pt]
\vspace{3pt}
\begin{tcolorbox}[colback=ofbTeal!8,colframe=ofbTeal,boxrule=0.6pt,arc=2pt,left=4pt,right=4pt,top=3pt,bottom=3pt]
{\small\bfseries Key line: }{\small Surprisingly, the phase jump does not break the EKF.}
\end{tcolorbox}
\end{minipage}
\clearpage
\noindent\begin{minipage}[t]{0.585\textwidth}\vspace{0pt}
{\setlength{\fboxsep}{0pt}\setlength{\fboxrule}{0.4pt}\fbox{\includegraphics[width=\linewidth]{build/guion/g-26.png}}}
\end{minipage}\hfill
\begin{minipage}[t]{0.395\textwidth}\vspace{2pt}
{\footnotesize\color{ofbMuted}Slide 26 / 60 \quad $\sim$25s \quad cum 10:56}\hfill{\footnotesize\bfseries\color{ofbTeal}CORE}\\[2pt]
{\large\bfseries\color{ofbTeal} Composite IBR sequence: hardest case}\\[6pt]
{\small\itshape\color{ofbMuted}Say:}\\[2pt]
\begin{itemize}[leftmargin=1.1em,itemsep=3pt,topsep=2pt]
\item Scenario E stacks it all: harmonics, two phase jumps, a RoCoF segment, ringdown, and impulsive noise.
\item Spectral leakage, loop re-lock and model mismatch all hit at the same time.
\end{itemize}
\vspace{2pt}{\small\color{ofbBlue}$\rightarrow$ \itshape And here the leaderboard falls apart.}\\[2pt]
\vspace{3pt}
\begin{tcolorbox}[colback=ofbTeal!8,colframe=ofbTeal,boxrule=0.6pt,arc=2pt,left=4pt,right=4pt,top=3pt,bottom=3pt]
{\small\bfseries Key line: }{\small This is the hardest and most realistic case in the suite.}
\end{tcolorbox}
\end{minipage}
\clearpage
\noindent\begin{minipage}[t]{0.585\textwidth}\vspace{0pt}
{\setlength{\fboxsep}{0pt}\setlength{\fboxrule}{0.4pt}\fbox{\includegraphics[width=\linewidth]{build/guion/g-27.png}}}
\end{minipage}\hfill
\begin{minipage}[t]{0.395\textwidth}\vspace{2pt}
{\footnotesize\color{ofbMuted}Slide 27 / 60 \quad $\sim$35s \quad cum 11:31}\hfill{\footnotesize\bfseries\color{ofbTeal}CORE}\\[2pt]
{\large\bfseries\color{ofbTeal} Scenario E: the leaderboard breaks}\\[6pt]
{\small\itshape\color{ofbMuted}Say:}\\[2pt]
\begin{itemize}[leftmargin=1.1em,itemsep=3pt,topsep=2pt]
\item Every usable method collapses to about one-point-one hertz; ZCD diverges to four hundred hertz.
\item And the rankings by RMSE, by trip-risk, and by CPU all disagree with each other.
\end{itemize}
\vspace{2pt}{\small\color{ofbBlue}$\rightarrow$ \itshape Let me make the trip-risk part visual.}\\[2pt]
\vspace{3pt}
\begin{tcolorbox}[colback=ofbTeal!8,colframe=ofbTeal,boxrule=0.6pt,arc=2pt,left=4pt,right=4pt,top=3pt,bottom=3pt]
{\small\bfseries Key line: }{\small Under stacked stress there is simply no accurate winner.}
\end{tcolorbox}
\end{minipage}
\clearpage
\noindent\begin{minipage}[t]{0.585\textwidth}\vspace{0pt}
{\setlength{\fboxsep}{0pt}\setlength{\fboxrule}{0.4pt}\fbox{\includegraphics[width=\linewidth]{build/guion/g-28.png}}}
\end{minipage}\hfill
\begin{minipage}[t]{0.395\textwidth}\vspace{2pt}
{\footnotesize\color{ofbMuted}Slide 28 / 60 \quad $\sim$30s \quad cum 12:01}\hfill{\footnotesize\bfseries\color{ofbBlue}FIG}\\[2pt]
{\large\bfseries\color{ofbBlue} Trip-risk evidence (figure)}\\[6pt]
{\small\itshape\color{ofbMuted}Say:}\\[2pt]
\begin{itemize}[leftmargin=1.1em,itemsep=3pt,topsep=2pt]
\item In Scenario D the state-space filters hold zero exposure while windowed methods pass 1.2 seconds.
\item In Scenario E, everyone carries exposure.
\end{itemize}
\vspace{2pt}{\small\color{ofbBlue}$\rightarrow$ \itshape And you can watch the divergence happen in real time.}\\[2pt]
\vspace{3pt}
\begin{tcolorbox}[colback=ofbBlue!8,colframe=ofbBlue,boxrule=0.6pt,arc=2pt,left=4pt,right=4pt,top=3pt,bottom=3pt]
{\small\bfseries Key line: }{\small Trip-risk is a different ranking from RMSE.}
\end{tcolorbox}
\end{minipage}
\clearpage
\noindent\begin{minipage}[t]{0.585\textwidth}\vspace{0pt}
{\setlength{\fboxsep}{0pt}\setlength{\fboxrule}{0.4pt}\fbox{\includegraphics[width=\linewidth]{build/guion/g-29.png}}}
\end{minipage}\hfill
\begin{minipage}[t]{0.395\textwidth}\vspace{2pt}
{\footnotesize\color{ofbMuted}Slide 29 / 60 \quad $\sim$30s \quad cum 12:31}\hfill{\footnotesize\bfseries\color{ofbBlue}FIG}\\[2pt]
{\large\bfseries\color{ofbBlue} Event responses (figure)}\\[6pt]
{\small\itshape\color{ofbMuted}Say:}\\[2pt]
\begin{itemize}[leftmargin=1.1em,itemsep=3pt,topsep=2pt]
\item True frequency in black, and each estimator trying to track it.
\item You can literally see EKF and IpDFT run away on the ramp, while RA-EKF stays locked.
\end{itemize}
\vspace{2pt}{\small\color{ofbBlue}$\rightarrow$ \itshape So what does all of this mean if you have to choose?}\\[2pt]
\vspace{3pt}
\begin{tcolorbox}[colback=ofbBlue!8,colframe=ofbBlue,boxrule=0.6pt,arc=2pt,left=4pt,right=4pt,top=3pt,bottom=3pt]
{\small\bfseries Key line: }{\small This is the intermittent divergence, seen with your own eyes.}
\end{tcolorbox}
\end{minipage}
\clearpage
\noindent\begin{minipage}[t]{0.585\textwidth}\vspace{0pt}
{\setlength{\fboxsep}{0pt}\setlength{\fboxrule}{0.4pt}\fbox{\includegraphics[width=\linewidth]{build/guion/g-30.png}}}
\end{minipage}\hfill
\begin{minipage}[t]{0.395\textwidth}\vspace{2pt}
{\footnotesize\color{ofbMuted}Slide 30 / 60 \quad $\sim$30s \quad cum 13:01}\hfill{\footnotesize\bfseries\color{ofbTeal}CORE}\\[2pt]
{\large\bfseries\color{ofbTeal} Result comparison: winners by question}\\[6pt]
{\small\itshape\color{ofbMuted}Say:}\\[2pt]
\begin{itemize}[leftmargin=1.1em,itemsep=3pt,topsep=2pt]
\item Ask for clean accuracy and you get the state-space and SOGI-FLL.
\item Ask for ramp robustness and you get RA-EKF. Ask about multi-event and there is no winner. Ask about feasibility and the expensive methods drop out.
\end{itemize}
\vspace{2pt}{\small\color{ofbBlue}$\rightarrow$ \itshape Feasibility brings us to cost.}\\[2pt]
\vspace{3pt}
\begin{tcolorbox}[colback=ofbTeal!8,colframe=ofbTeal,boxrule=0.6pt,arc=2pt,left=4pt,right=4pt,top=3pt,bottom=3pt]
{\small\bfseries Key line: }{\small Different question, different winner, every time.}
\end{tcolorbox}
\end{minipage}
\clearpage
\noindent\begin{minipage}[t]{0.585\textwidth}\vspace{0pt}
{\setlength{\fboxsep}{0pt}\setlength{\fboxrule}{0.4pt}\fbox{\includegraphics[width=\linewidth]{build/guion/g-31.png}}}
\end{minipage}\hfill
\begin{minipage}[t]{0.395\textwidth}\vspace{2pt}
{\footnotesize\color{ofbMuted}Slide 31 / 60 \quad $\sim$28s \quad cum 13:29}\hfill{\footnotesize\bfseries\color{ofbTeal}CORE}\\[2pt]
{\large\bfseries\color{ofbTeal} Computational feasibility}\\[6pt]
{\small\itshape\color{ofbMuted}Say:}\\[2pt]
\begin{itemize}[leftmargin=1.1em,itemsep=3pt,topsep=2pt]
\item The loop and state-space methods run in nine to seventeen microseconds per sample, which is embedded-feasible.
\item The spectral and data-driven methods are a hundred to a thousand times slower.
\end{itemize}
\vspace{2pt}{\small\color{ofbBlue}$\rightarrow$ \itshape Put cost and risk together and the picture is clear.}\\[2pt]
\vspace{3pt}
\begin{tcolorbox}[colback=ofbTeal!8,colframe=ofbTeal,boxrule=0.6pt,arc=2pt,left=4pt,right=4pt,top=3pt,bottom=3pt]
{\small\bfseries Key line: }{\small Cost alone already removes several methods from a relay.}
\end{tcolorbox}
\end{minipage}
\clearpage
\noindent\begin{minipage}[t]{0.585\textwidth}\vspace{0pt}
{\setlength{\fboxsep}{0pt}\setlength{\fboxrule}{0.4pt}\fbox{\includegraphics[width=\linewidth]{build/guion/g-32.png}}}
\end{minipage}\hfill
\begin{minipage}[t]{0.395\textwidth}\vspace{2pt}
{\footnotesize\color{ofbMuted}Slide 32 / 60 \quad $\sim$25s \quad cum 13:54}\hfill{\footnotesize\bfseries\color{ofbBlue}FIG}\\[2pt]
{\large\bfseries\color{ofbBlue} Deployment limits (figure)}\\[6pt]
{\small\itshape\color{ofbMuted}Say:}\\[2pt]
\begin{itemize}[leftmargin=1.1em,itemsep=3pt,topsep=2pt]
\item Deployment is the combination of accuracy, trip exposure, cost, and divergence-robustness.
\item And the low-cost cluster, the state-space and loop methods, is also the low-trip cluster.
\end{itemize}
\vspace{2pt}{\small\color{ofbBlue}$\rightarrow$ \itshape We can summarize each method's balance in one shape.}\\[2pt]
\vspace{3pt}
\begin{tcolorbox}[colback=ofbBlue!8,colframe=ofbBlue,boxrule=0.6pt,arc=2pt,left=4pt,right=4pt,top=3pt,bottom=3pt]
{\small\bfseries Key line: }{\small The practical winners sit together at low cost and low trip.}
\end{tcolorbox}
\end{minipage}
\clearpage
\noindent\begin{minipage}[t]{0.585\textwidth}\vspace{0pt}
{\setlength{\fboxsep}{0pt}\setlength{\fboxrule}{0.4pt}\fbox{\includegraphics[width=\linewidth]{build/guion/g-33.png}}}
\end{minipage}\hfill
\begin{minipage}[t]{0.395\textwidth}\vspace{2pt}
{\footnotesize\color{ofbMuted}Slide 33 / 60 \quad $\sim$25s \quad cum 14:19}\hfill{\footnotesize\bfseries\color{ofbBlue}FIG}\\[2pt]
{\large\bfseries\color{ofbBlue} Balance profile (radar)}\\[6pt]
{\small\itshape\color{ofbMuted}Say:}\\[2pt]
\begin{itemize}[leftmargin=1.1em,itemsep=3pt,topsep=2pt]
\item Each axis is one scenario, and bigger means more robust.
\item RA-EKF and SOGI-FLL stay robust on every axis; EKF and IpDFT collapse on the ramp axis.
\end{itemize}
\vspace{2pt}{\small\color{ofbBlue}$\rightarrow$ \itshape Now, the most important slide about compliance testing.}\\[2pt]
\vspace{3pt}
\begin{tcolorbox}[colback=ofbBlue!8,colframe=ofbBlue,boxrule=0.6pt,arc=2pt,left=4pt,right=4pt,top=3pt,bottom=3pt]
{\small\bfseries Key line: }{\small No method is strong on every axis.}
\end{tcolorbox}
\end{minipage}
\clearpage
\noindent\begin{minipage}[t]{0.585\textwidth}\vspace{0pt}
{\setlength{\fboxsep}{0pt}\setlength{\fboxrule}{0.4pt}\fbox{\includegraphics[width=\linewidth]{build/guion/g-34.png}}}
\end{minipage}\hfill
\begin{minipage}[t]{0.395\textwidth}\vspace{2pt}
{\footnotesize\color{ofbMuted}Slide 34 / 60 \quad $\sim$35s \quad cum 14:54}\hfill{\footnotesize\bfseries\color{ofbTeal}CORE}\\[2pt]
{\large\bfseries\color{ofbTeal} Compliance heatmap (build-up)}\\[6pt]
{\small\itshape\color{ofbMuted}Say:}\\[2pt]
\begin{itemize}[leftmargin=1.1em,itemsep=3pt,topsep=2pt]
\item Look at the standard columns first: almost everyone passes step, ramp, and modulation.
\item CLICK. But the same methods fail on phase jump and multi-event.
\end{itemize}
\vspace{2pt}{\small\color{ofbBlue}$\rightarrow$ \itshape Let me distill that into our empirical claims.}\\[2pt]
\vspace{3pt}
\begin{tcolorbox}[colback=ofbTeal!8,colframe=ofbTeal,boxrule=0.6pt,arc=2pt,left=4pt,right=4pt,top=3pt,bottom=3pt]
{\small\bfseries Key line: }{\small Passing the standard tests does not predict event readiness.}
\end{tcolorbox}
\end{minipage}
\clearpage
\noindent\begin{minipage}[t]{0.585\textwidth}\vspace{0pt}
{\setlength{\fboxsep}{0pt}\setlength{\fboxrule}{0.4pt}\fbox{\includegraphics[width=\linewidth]{build/guion/g-35.png}}}
\end{minipage}\hfill
\begin{minipage}[t]{0.395\textwidth}\vspace{2pt}
{\footnotesize\color{ofbMuted}Slide 35 / 60 \quad $\sim$30s \quad cum 15:24}\hfill{\footnotesize\bfseries\color{ofbTeal}CORE}\\[2pt]
{\large\bfseries\color{ofbTeal} Empirical claims}\\[6pt]
{\small\itshape\color{ofbMuted}Say:}\\[2pt]
\begin{itemize}[leftmargin=1.1em,itemsep=3pt,topsep=2pt]
\item Four claims. Divergence, not phase jumps, is the dominant failure mode.
\item Pass-fail leaves event readiness unresolved; the three metrics pick different winners; and feasibility changes the ranking.
\end{itemize}
\vspace{2pt}{\small\color{ofbBlue}$\rightarrow$ \itshape Now, does this pattern survive when we scale up?}\\[2pt]
\vspace{3pt}
\begin{tcolorbox}[colback=ofbTeal!8,colframe=ofbTeal,boxrule=0.6pt,arc=2pt,left=4pt,right=4pt,top=3pt,bottom=3pt]
{\small\bfseries Key line: }{\small Four claims, one theme: always report the conditions.}
\end{tcolorbox}
\end{minipage}
\clearpage
\noindent\begin{minipage}[t]{0.585\textwidth}\vspace{0pt}
{\setlength{\fboxsep}{0pt}\setlength{\fboxrule}{0.4pt}\fbox{\includegraphics[width=\linewidth]{build/guion/g-36.png}}}
\end{minipage}\hfill
\begin{minipage}[t]{0.395\textwidth}\vspace{2pt}
{\footnotesize\color{ofbMuted}Slide 36 / 60 \quad $\sim$6s \quad cum 15:30}\hfill{\footnotesize\bfseries\color{ofbMuted}QUICK}\\[2pt]
{\large\bfseries\color{ofbMuted} Expanded and ATLAS (section)}\\[6pt]
\vspace{3pt}
\begin{tcolorbox}[colback=ofbMuted!8,colframe=ofbMuted,boxrule=0.6pt,arc=2pt,left=4pt,right=4pt,top=3pt,bottom=3pt]
{\small\bfseries Key line: }{\small Does the pattern hold across families and severity? That is ATLAS.}
\end{tcolorbox}
\end{minipage}
\clearpage
\noindent\begin{minipage}[t]{0.585\textwidth}\vspace{0pt}
{\setlength{\fboxsep}{0pt}\setlength{\fboxrule}{0.4pt}\fbox{\includegraphics[width=\linewidth]{build/guion/g-37.png}}}
\end{minipage}\hfill
\begin{minipage}[t]{0.395\textwidth}\vspace{2pt}
{\footnotesize\color{ofbMuted}Slide 37 / 60 \quad $\sim$22s \quad cum 15:52}\hfill{\footnotesize\bfseries\color{ofbTeal}CORE}\\[2pt]
{\large\bfseries\color{ofbTeal} Expanded stress analysis}\\[6pt]
{\small\itshape\color{ofbMuted}Say:}\\[2pt]
\begin{itemize}[leftmargin=1.1em,itemsep=3pt,topsep=2pt]
\item We scale up to a full Monte-Carlo benchmark: eighteen estimators, thirty-three scenarios, five seeds.
\item And then the ATLAS severity sweeps on top.
\end{itemize}
\vspace{2pt}{\small\color{ofbBlue}$\rightarrow$ \itshape First, do the family winners stay fixed?}\\[2pt]
\vspace{3pt}
\begin{tcolorbox}[colback=ofbTeal!8,colframe=ofbTeal,boxrule=0.6pt,arc=2pt,left=4pt,right=4pt,top=3pt,bottom=3pt]
{\small\bfseries Key line: }{\small We scale the test, we do not just repeat it.}
\end{tcolorbox}
\end{minipage}
\clearpage
\noindent\begin{minipage}[t]{0.585\textwidth}\vspace{0pt}
{\setlength{\fboxsep}{0pt}\setlength{\fboxrule}{0.4pt}\fbox{\includegraphics[width=\linewidth]{build/guion/g-38.png}}}
\end{minipage}\hfill
\begin{minipage}[t]{0.395\textwidth}\vspace{2pt}
{\footnotesize\color{ofbMuted}Slide 38 / 60 \quad $\sim$22s \quad cum 16:14}\hfill{\footnotesize\bfseries\color{ofbTeal}CORE}\\[2pt]
{\large\bfseries\color{ofbTeal} Expanded benchmark: family winners}\\[6pt]
{\small\itshape\color{ofbMuted}Say:}\\[2pt]
\begin{itemize}[leftmargin=1.1em,itemsep=3pt,topsep=2pt]
\item They do not. The winner rotates: one method leads the step, another the ramp, another harmonics, another the multi-event.
\end{itemize}
\vspace{2pt}{\small\color{ofbBlue}$\rightarrow$ \itshape ATLAS lets us see why, by sweeping severity.}\\[2pt]
\vspace{3pt}
\begin{tcolorbox}[colback=ofbTeal!8,colframe=ofbTeal,boxrule=0.6pt,arc=2pt,left=4pt,right=4pt,top=3pt,bottom=3pt]
{\small\bfseries Key line: }{\small The leaderboard is regime-dependent, never absolute.}
\end{tcolorbox}
\end{minipage}
\clearpage
\noindent\begin{minipage}[t]{0.585\textwidth}\vspace{0pt}
{\setlength{\fboxsep}{0pt}\setlength{\fboxrule}{0.4pt}\fbox{\includegraphics[width=\linewidth]{build/guion/g-39.png}}}
\end{minipage}\hfill
\begin{minipage}[t]{0.395\textwidth}\vspace{2pt}
{\footnotesize\color{ofbMuted}Slide 39 / 60 \quad $\sim$25s \quad cum 16:39}\hfill{\footnotesize\bfseries\color{ofbTeal}CORE}\\[2pt]
{\large\bfseries\color{ofbTeal} ATLAS severity analysis}\\[6pt]
{\small\itshape\color{ofbMuted}Say:}\\[2pt]
\begin{itemize}[leftmargin=1.1em,itemsep=3pt,topsep=2pt]
\item ATLAS sweeps each stress by severity and by sign, under a fixed policy, at thirty Monte-Carlo runs.
\item This is now paper-grade across all nine required sweeps.
\end{itemize}
\vspace{2pt}{\small\color{ofbBlue}$\rightarrow$ \itshape And the headline result is striking.}\\[2pt]
\vspace{3pt}
\begin{tcolorbox}[colback=ofbTeal!8,colframe=ofbTeal,boxrule=0.6pt,arc=2pt,left=4pt,right=4pt,top=3pt,bottom=3pt]
{\small\bfseries Key line: }{\small This is paper-grade severity evidence, not a pilot.}
\end{tcolorbox}
\end{minipage}
\clearpage
\noindent\begin{minipage}[t]{0.585\textwidth}\vspace{0pt}
{\setlength{\fboxsep}{0pt}\setlength{\fboxrule}{0.4pt}\fbox{\includegraphics[width=\linewidth]{build/guion/g-40.png}}}
\end{minipage}\hfill
\begin{minipage}[t]{0.395\textwidth}\vspace{2pt}
{\footnotesize\color{ofbMuted}Slide 40 / 60 \quad $\sim$30s \quad cum 17:09}\hfill{\footnotesize\bfseries\color{ofbTeal}CORE}\\[2pt]
{\large\bfseries\color{ofbTeal} ATLAS: metric choice changes the answer}\\[6pt]
{\small\itshape\color{ofbMuted}Say:}\\[2pt]
\begin{itemize}[leftmargin=1.1em,itemsep=3pt,topsep=2pt]
\item Across a hundred and thirteen severity levels, the preferred estimator changes when the objective changes.
\item In eighty-three of them, RMSE and trip-risk point to different methods.
\end{itemize}
\vspace{2pt}{\small\color{ofbBlue}$\rightarrow$ \itshape ATLAS also tells us where each method gives up.}\\[2pt]
\vspace{3pt}
\begin{tcolorbox}[colback=ofbTeal!8,colframe=ofbTeal,boxrule=0.6pt,arc=2pt,left=4pt,right=4pt,top=3pt,bottom=3pt]
{\small\bfseries Key line: }{\small The output is a selection map, not a leaderboard.}
\end{tcolorbox}
\end{minipage}
\clearpage
\noindent\begin{minipage}[t]{0.585\textwidth}\vspace{0pt}
{\setlength{\fboxsep}{0pt}\setlength{\fboxrule}{0.4pt}\fbox{\includegraphics[width=\linewidth]{build/guion/g-41.png}}}
\end{minipage}\hfill
\begin{minipage}[t]{0.395\textwidth}\vspace{2pt}
{\footnotesize\color{ofbMuted}Slide 41 / 60 \quad $\sim$22s \quad cum 17:31}\hfill{\footnotesize\bfseries\color{ofbBlue}FIG}\\[2pt]
{\large\bfseries\color{ofbBlue} ATLAS: where estimators stop being valid}\\[6pt]
{\small\itshape\color{ofbMuted}Say:}\\[2pt]
\begin{itemize}[leftmargin=1.1em,itemsep=3pt,topsep=2pt]
\item For each method, ATLAS finds the severity at which it crosses into failure.
\end{itemize}
\vspace{2pt}{\small\color{ofbBlue}$\rightarrow$ \itshape Let me show you a few of the most surprising sweeps.}\\[2pt]
\vspace{3pt}
\begin{tcolorbox}[colback=ofbBlue!8,colframe=ofbBlue,boxrule=0.6pt,arc=2pt,left=4pt,right=4pt,top=3pt,bottom=3pt]
{\small\bfseries Key line: }{\small That validity boundary is different for every family.}
\end{tcolorbox}
\end{minipage}
\clearpage
\noindent\begin{minipage}[t]{0.585\textwidth}\vspace{0pt}
{\setlength{\fboxsep}{0pt}\setlength{\fboxrule}{0.4pt}\fbox{\includegraphics[width=\linewidth]{build/guion/g-42.png}}}
\end{minipage}\hfill
\begin{minipage}[t]{0.395\textwidth}\vspace{2pt}
{\footnotesize\color{ofbMuted}Slide 42 / 60 \quad $\sim$25s \quad cum 17:56}\hfill{\footnotesize\bfseries\color{ofbTeal}CORE}\\[2pt]
{\large\bfseries\color{ofbTeal} AM, FM and interharmonics}\\[6pt]
{\small\itshape\color{ofbMuted}Say:}\\[2pt]
\begin{itemize}[leftmargin=1.1em,itemsep=3pt,topsep=2pt]
\item Modulation actually splits the story: AM is easy, and the model-based methods win.
\item FM and interharmonics are harder, and the spectral methods take over.
\end{itemize}
\vspace{2pt}{\small\color{ofbBlue}$\rightarrow$ \itshape Now the four findings I most want you to remember.}\\[2pt]
\vspace{3pt}
\begin{tcolorbox}[colback=ofbTeal!8,colframe=ofbTeal,boxrule=0.6pt,arc=2pt,left=4pt,right=4pt,top=3pt,bottom=3pt]
{\small\bfseries Key line: }{\small Even inside modulation, the winner switches.}
\end{tcolorbox}
\end{minipage}
\clearpage
\noindent\begin{minipage}[t]{0.585\textwidth}\vspace{0pt}
{\setlength{\fboxsep}{0pt}\setlength{\fboxrule}{0.4pt}\fbox{\includegraphics[width=\linewidth]{build/guion/g-43.png}}}
\end{minipage}\hfill
\begin{minipage}[t]{0.395\textwidth}\vspace{2pt}
{\footnotesize\color{ofbMuted}Slide 43 / 60 \quad $\sim$30s \quad cum 18:26}\hfill{\footnotesize\bfseries\color{ofbTeal}CORE}\\[2pt]
{\large\bfseries\color{ofbTeal} When the textbook filter diverges}\\[6pt]
{\small\itshape\color{ofbMuted}Say:}\\[2pt]
\begin{itemize}[leftmargin=1.1em,itemsep=3pt,topsep=2pt]
\item The textbook EKF tracks most seeds, then intermittently runs away to tens of hertz.
\item And it is not severity-driven: the worst seed actually had the gentlest ramp.
\end{itemize}
\vspace{2pt}{\small\color{ofbBlue}$\rightarrow$ \itshape The next finding is about direction.}\\[2pt]
\vspace{3pt}
\begin{tcolorbox}[colback=ofbTeal!8,colframe=ofbTeal,boxrule=0.6pt,arc=2pt,left=4pt,right=4pt,top=3pt,bottom=3pt]
{\small\bfseries Key line: }{\small The robustified RA-EKF removes the divergence on every seed.}
\end{tcolorbox}
\end{minipage}
\clearpage
\noindent\begin{minipage}[t]{0.585\textwidth}\vspace{0pt}
{\setlength{\fboxsep}{0pt}\setlength{\fboxrule}{0.4pt}\fbox{\includegraphics[width=\linewidth]{build/guion/g-44.png}}}
\end{minipage}\hfill
\begin{minipage}[t]{0.395\textwidth}\vspace{2pt}
{\footnotesize\color{ofbMuted}Slide 44 / 60 \quad $\sim$30s \quad cum 18:56}\hfill{\footnotesize\bfseries\color{ofbTeal}CORE}\\[2pt]
{\large\bfseries\color{ofbTeal} RoCoF has a sign}\\[6pt]
{\small\itshape\color{ofbMuted}Say:}\\[2pt]
\begin{itemize}[leftmargin=1.1em,itemsep=3pt,topsep=2pt]
\item RoCoF has a sign, and down-ramps are harder than up-ramps.
\item LKF is five-and-a-half times worse on a falling frequency; TKEO almost four times worse.
\end{itemize}
\vspace{2pt}{\small\color{ofbBlue}$\rightarrow$ \itshape Severity can also surprise us in the other direction.}\\[2pt]
\vspace{3pt}
\begin{tcolorbox}[colback=ofbTeal!8,colframe=ofbTeal,boxrule=0.6pt,arc=2pt,left=4pt,right=4pt,top=3pt,bottom=3pt]
{\small\bfseries Key line: }{\small A rising-only test can certify a method that fails when frequency falls, the protection-critical case.}
\end{tcolorbox}
\end{minipage}
\clearpage
\noindent\begin{minipage}[t]{0.585\textwidth}\vspace{0pt}
{\setlength{\fboxsep}{0pt}\setlength{\fboxrule}{0.4pt}\fbox{\includegraphics[width=\linewidth]{build/guion/g-45.png}}}
\end{minipage}\hfill
\begin{minipage}[t]{0.395\textwidth}\vspace{2pt}
{\footnotesize\color{ofbMuted}Slide 45 / 60 \quad $\sim$28s \quad cum 19:24}\hfill{\footnotesize\bfseries\color{ofbTeal}CORE}\\[2pt]
{\large\bfseries\color{ofbTeal} Phase-jump severity is non-monotone}\\[6pt]
{\small\itshape\color{ofbMuted}Say:}\\[2pt]
\begin{itemize}[leftmargin=1.1em,itemsep=3pt,topsep=2pt]
\item Phase-jump severity is non-monotone: the worst case is not a hundred and eighty degrees.
\item Most methods peak around ninety to a hundred and twenty degrees and then recover, because a half-cycle jump partly cancels itself.
\end{itemize}
\vspace{2pt}{\small\color{ofbBlue}$\rightarrow$ \itshape Severity also reshuffles the winners completely.}\\[2pt]
\vspace{3pt}
\begin{tcolorbox}[colback=ofbTeal!8,colframe=ofbTeal,boxrule=0.6pt,arc=2pt,left=4pt,right=4pt,top=3pt,bottom=3pt]
{\small\bfseries Key line: }{\small Bigger is not always worse, so you must sweep the range.}
\end{tcolorbox}
\end{minipage}
\clearpage
\noindent\begin{minipage}[t]{0.585\textwidth}\vspace{0pt}
{\setlength{\fboxsep}{0pt}\setlength{\fboxrule}{0.4pt}\fbox{\includegraphics[width=\linewidth]{build/guion/g-46.png}}}
\end{minipage}\hfill
\begin{minipage}[t]{0.395\textwidth}\vspace{2pt}
{\footnotesize\color{ofbMuted}Slide 46 / 60 \quad $\sim$28s \quad cum 19:52}\hfill{\footnotesize\bfseries\color{ofbBlue}FIG}\\[2pt]
{\large\bfseries\color{ofbBlue} No universal winner (map)}\\[6pt]
{\small\itshape\color{ofbMuted}Say:}\\[2pt]
\begin{itemize}[leftmargin=1.1em,itemsep=3pt,topsep=2pt]
\item This map says everything: the champion rotates by RoCoF regime.
\item Koopman at low rates, ESPRIT in the middle, SOGI-PLL when it gets severe.
\end{itemize}
\vspace{2pt}{\small\color{ofbBlue}$\rightarrow$ \itshape And one estimator manages to be both the best and the worst.}\\[2pt]
\vspace{3pt}
\begin{tcolorbox}[colback=ofbBlue!8,colframe=ofbBlue,boxrule=0.6pt,arc=2pt,left=4pt,right=4pt,top=3pt,bottom=3pt]
{\small\bfseries Key line: }{\small No universal winner; the deliverable is the map itself.}
\end{tcolorbox}
\end{minipage}
\clearpage
\noindent\begin{minipage}[t]{0.585\textwidth}\vspace{0pt}
{\setlength{\fboxsep}{0pt}\setlength{\fboxrule}{0.4pt}\fbox{\includegraphics[width=\linewidth]{build/guion/g-47.png}}}
\end{minipage}\hfill
\begin{minipage}[t]{0.395\textwidth}\vspace{2pt}
{\footnotesize\color{ofbMuted}Slide 47 / 60 \quad $\sim$40s \quad cum 20:32}\hfill{\footnotesize\bfseries\color{ofbTeal}CORE}\\[2pt]
{\large\bfseries\color{ofbTeal} The ZCD paradox (build-up)}\\[6pt]
{\small\itshape\color{ofbMuted}Say:}\\[2pt]
\begin{itemize}[leftmargin=1.1em,itemsep=3pt,topsep=2pt]
\item Below fifteen percent harmonic distortion, zero-crossing detection is the most accurate method we have.
\item CLICK. Add broadband noise and it explodes to three thousand hertz.
\item CLICK. Push past twenty percent distortion and it collapses to tens of hertz.
\end{itemize}
\vspace{2pt}{\small\color{ofbBlue}$\rightarrow$ \itshape The same asymmetry shows up in voltage events.}\\[2pt]
\vspace{3pt}
\begin{tcolorbox}[colback=ofbTeal!8,colframe=ofbTeal,boxrule=0.6pt,arc=2pt,left=4pt,right=4pt,top=3pt,bottom=3pt]
{\small\bfseries Key line: }{\small The same estimator is the best and the worst, depending on the stress.}
\end{tcolorbox}
\end{minipage}
\clearpage
\noindent\begin{minipage}[t]{0.585\textwidth}\vspace{0pt}
{\setlength{\fboxsep}{0pt}\setlength{\fboxrule}{0.4pt}\fbox{\includegraphics[width=\linewidth]{build/guion/g-48.png}}}
\end{minipage}\hfill
\begin{minipage}[t]{0.395\textwidth}\vspace{2pt}
{\footnotesize\color{ofbMuted}Slide 48 / 60 \quad $\sim$25s \quad cum 20:57}\hfill{\footnotesize\bfseries\color{ofbTeal}CORE}\\[2pt]
{\large\bfseries\color{ofbTeal} Sags and swells are not symmetric}\\[6pt]
{\small\itshape\color{ofbMuted}Say:}\\[2pt]
\begin{itemize}[leftmargin=1.1em,itemsep=3pt,topsep=2pt]
\item Voltage magnitude is not symmetric either: a deep sag is far harder than a swell.
\item ZCD goes from millihertz on a swell to hundreds of hertz on a deep sag.
\end{itemize}
\vspace{2pt}{\small\color{ofbBlue}$\rightarrow$ \itshape And there is a clean physical reason.}\\[2pt]
\vspace{3pt}
\begin{tcolorbox}[colback=ofbTeal!8,colframe=ofbTeal,boxrule=0.6pt,arc=2pt,left=4pt,right=4pt,top=3pt,bottom=3pt]
{\small\bfseries Key line: }{\small Direction matters; you have to sweep both signs.}
\end{tcolorbox}
\end{minipage}
\clearpage
\noindent\begin{minipage}[t]{0.585\textwidth}\vspace{0pt}
{\setlength{\fboxsep}{0pt}\setlength{\fboxrule}{0.4pt}\fbox{\includegraphics[width=\linewidth]{build/guion/g-49.png}}}
\end{minipage}\hfill
\begin{minipage}[t]{0.395\textwidth}\vspace{2pt}
{\footnotesize\color{ofbMuted}Slide 49 / 60 \quad $\sim$25s \quad cum 21:22}\hfill{\footnotesize\bfseries\color{ofbBlue}FIG}\\[2pt]
{\large\bfseries\color{ofbBlue} Why a sag breaks zero-crossing}\\[6pt]
{\small\itshape\color{ofbMuted}Say:}\\[2pt]
\begin{itemize}[leftmargin=1.1em,itemsep=3pt,topsep=2pt]
\item The timing error of a zero-crossing scales with noise over amplitude times frequency.
\item A deep sag flattens the slope right at the crossing, so noise creates false crossings.
\end{itemize}
\vspace{2pt}{\small\color{ofbBlue}$\rightarrow$ \itshape Harmonics tell a similar family-by-family story.}\\[2pt]
\vspace{3pt}
\begin{tcolorbox}[colback=ofbBlue!8,colframe=ofbBlue,boxrule=0.6pt,arc=2pt,left=4pt,right=4pt,top=3pt,bottom=3pt]
{\small\bfseries Key line: }{\small This is physics, not a software bug.}
\end{tcolorbox}
\end{minipage}
\clearpage
\noindent\begin{minipage}[t]{0.585\textwidth}\vspace{0pt}
{\setlength{\fboxsep}{0pt}\setlength{\fboxrule}{0.4pt}\fbox{\includegraphics[width=\linewidth]{build/guion/g-50.png}}}
\end{minipage}\hfill
\begin{minipage}[t]{0.395\textwidth}\vspace{2pt}
{\footnotesize\color{ofbMuted}Slide 50 / 60 \quad $\sim$25s \quad cum 21:47}\hfill{\footnotesize\bfseries\color{ofbTeal}CORE}\\[2pt]
{\large\bfseries\color{ofbTeal} Harmonic degradation by family}\\[6pt]
{\small\itshape\color{ofbMuted}Say:}\\[2pt]
\begin{itemize}[leftmargin=1.1em,itemsep=3pt,topsep=2pt]
\item Harmonic degradation depends on the family: loop methods like ZCD win at low distortion but collapse above fifteen percent.
\item Window methods degrade gently and end up leading at extreme distortion.
\end{itemize}
\vspace{2pt}{\small\color{ofbBlue}$\rightarrow$ \itshape And cost is the last dimension.}\\[2pt]
\vspace{3pt}
\begin{tcolorbox}[colback=ofbTeal!8,colframe=ofbTeal,boxrule=0.6pt,arc=2pt,left=4pt,right=4pt,top=3pt,bottom=3pt]
{\small\bfseries Key line: }{\small Low-distortion accuracy does not predict high-distortion survival.}
\end{tcolorbox}
\end{minipage}
\clearpage
\noindent\begin{minipage}[t]{0.585\textwidth}\vspace{0pt}
{\setlength{\fboxsep}{0pt}\setlength{\fboxrule}{0.4pt}\fbox{\includegraphics[width=\linewidth]{build/guion/g-51.png}}}
\end{minipage}\hfill
\begin{minipage}[t]{0.395\textwidth}\vspace{2pt}
{\footnotesize\color{ofbMuted}Slide 51 / 60 \quad $\sim$22s \quad cum 22:09}\hfill{\footnotesize\bfseries\color{ofbBlue}FIG}\\[2pt]
{\large\bfseries\color{ofbBlue} CPU-accuracy Pareto}\\[6pt]
{\small\itshape\color{ofbMuted}Say:}\\[2pt]
\begin{itemize}[leftmargin=1.1em,itemsep=3pt,topsep=2pt]
\item On the cost-accuracy Pareto, the spectral and data-driven methods buy a little accuracy for a hundred to a thousand times more compute.
\end{itemize}
\vspace{2pt}{\small\color{ofbBlue}$\rightarrow$ \itshape So, what do we actually deploy?}\\[2pt]
\vspace{3pt}
\begin{tcolorbox}[colback=ofbBlue!8,colframe=ofbBlue,boxrule=0.6pt,arc=2pt,left=4pt,right=4pt,top=3pt,bottom=3pt]
{\small\bfseries Key line: }{\small For relay-class timing, that trade is usually not worth it.}
\end{tcolorbox}
\end{minipage}
\clearpage
\noindent\begin{minipage}[t]{0.585\textwidth}\vspace{0pt}
{\setlength{\fboxsep}{0pt}\setlength{\fboxrule}{0.4pt}\fbox{\includegraphics[width=\linewidth]{build/guion/g-52.png}}}
\end{minipage}\hfill
\begin{minipage}[t]{0.395\textwidth}\vspace{2pt}
{\footnotesize\color{ofbMuted}Slide 52 / 60 \quad $\sim$6s \quad cum 22:15}\hfill{\footnotesize\bfseries\color{ofbMuted}QUICK}\\[2pt]
{\large\bfseries\color{ofbMuted} Findings and deployment (section)}\\[6pt]
\vspace{3pt}
\begin{tcolorbox}[colback=ofbMuted!8,colframe=ofbMuted,boxrule=0.6pt,arc=2pt,left=4pt,right=4pt,top=3pt,bottom=3pt]
{\small\bfseries Key line: }{\small So what do we actually deploy?}
\end{tcolorbox}
\end{minipage}
\clearpage
\noindent\begin{minipage}[t]{0.585\textwidth}\vspace{0pt}
{\setlength{\fboxsep}{0pt}\setlength{\fboxrule}{0.4pt}\fbox{\includegraphics[width=\linewidth]{build/guion/g-53.png}}}
\end{minipage}\hfill
\begin{minipage}[t]{0.395\textwidth}\vspace{2pt}
{\footnotesize\color{ofbMuted}Slide 53 / 60 \quad $\sim$35s \quad cum 22:50}\hfill{\footnotesize\bfseries\color{ofbTeal}CORE}\\[2pt]
{\large\bfseries\color{ofbTeal} Three findings}\\[6pt]
{\small\itshape\color{ofbMuted}Say:}\\[2pt]
\begin{itemize}[leftmargin=1.1em,itemsep=3pt,topsep=2pt]
\item Three findings. One, there is a latency-versus-stress trade-off.
\item Two, robustness beats point accuracy: RA-EKF's explicit RoCoF state keeps it stable exactly where EKF and IpDFT diverge.
\item Three, compliance tests are not enough for inverter events.
\end{itemize}
\vspace{2pt}{\small\color{ofbBlue}$\rightarrow$ \itshape Let me be honest about the limits.}\\[2pt]
\vspace{3pt}
\begin{tcolorbox}[colback=ofbTeal!8,colframe=ofbTeal,boxrule=0.6pt,arc=2pt,left=4pt,right=4pt,top=3pt,bottom=3pt]
{\small\bfseries Key line: }{\small Robustness, not accuracy, is the deployment discriminator.}
\end{tcolorbox}
\end{minipage}
\clearpage
\noindent\begin{minipage}[t]{0.585\textwidth}\vspace{0pt}
{\setlength{\fboxsep}{0pt}\setlength{\fboxrule}{0.4pt}\fbox{\includegraphics[width=\linewidth]{build/guion/g-54.png}}}
\end{minipage}\hfill
\begin{minipage}[t]{0.395\textwidth}\vspace{2pt}
{\footnotesize\color{ofbMuted}Slide 54 / 60 \quad $\sim$28s \quad cum 23:18}\hfill{\footnotesize\bfseries\color{ofbTeal}CORE}\\[2pt]
{\large\bfseries\color{ofbTeal} Validity limits and next steps}\\[6pt]
{\small\itshape\color{ofbMuted}Say:}\\[2pt]
\begin{itemize}[leftmargin=1.1em,itemsep=3pt,topsep=2pt]
\item On validity: the core means are five seeds, and the ATLAS sweeps are paper-grade at thirty.
\item The next step is journal-grade at a hundred, plus the subspace and machine-learning methods.
\end{itemize}
\vspace{2pt}{\small\color{ofbBlue}$\rightarrow$ \itshape From here it becomes a recommendation.}\\[2pt]
\vspace{3pt}
\begin{tcolorbox}[colback=ofbTeal!8,colframe=ofbTeal,boxrule=0.6pt,arc=2pt,left=4pt,right=4pt,top=3pt,bottom=3pt]
{\small\bfseries Key line: }{\small We are explicit about what is diagnostic and what is paper-grade.}
\end{tcolorbox}
\end{minipage}
\clearpage
\noindent\begin{minipage}[t]{0.585\textwidth}\vspace{0pt}
{\setlength{\fboxsep}{0pt}\setlength{\fboxrule}{0.4pt}\fbox{\includegraphics[width=\linewidth]{build/guion/g-55.png}}}
\end{minipage}\hfill
\begin{minipage}[t]{0.395\textwidth}\vspace{2pt}
{\footnotesize\color{ofbMuted}Slide 55 / 60 \quad $\sim$30s \quad cum 23:48}\hfill{\footnotesize\bfseries\color{ofbTeal}CORE}\\[2pt]
{\large\bfseries\color{ofbTeal} Recommendation map}\\[6pt]
{\small\itshape\color{ofbMuted}Say:}\\[2pt]
\begin{itemize}[leftmargin=1.1em,itemsep=3pt,topsep=2pt]
\item For fast frequency response and anti-islanding, RA-EKF is the strongest benchmark-backed choice.
\item PLL and EKF are solid low-cost baselines, and IpDFT or TFT work where some observation delay is acceptable.
\end{itemize}
\vspace{2pt}{\small\color{ofbBlue}$\rightarrow$ \itshape And we wrap that into a qualification idea.}\\[2pt]
\vspace{3pt}
\begin{tcolorbox}[colback=ofbTeal!8,colframe=ofbTeal,boxrule=0.6pt,arc=2pt,left=4pt,right=4pt,top=3pt,bottom=3pt]
{\small\bfseries Key line: }{\small We deliver a recommendation map, not a single ranking.}
\end{tcolorbox}
\end{minipage}
\clearpage
\noindent\begin{minipage}[t]{0.585\textwidth}\vspace{0pt}
{\setlength{\fboxsep}{0pt}\setlength{\fboxrule}{0.4pt}\fbox{\includegraphics[width=\linewidth]{build/guion/g-56.png}}}
\end{minipage}\hfill
\begin{minipage}[t]{0.395\textwidth}\vspace{2pt}
{\footnotesize\color{ofbMuted}Slide 56 / 60 \quad $\sim$22s \quad cum 24:10}\hfill{\footnotesize\bfseries\color{ofbBlue}FIG}\\[2pt]
{\large\bfseries\color{ofbBlue} Qualification profile}\\[6pt]
{\small\itshape\color{ofbMuted}Say:}\\[2pt]
\begin{itemize}[leftmargin=1.1em,itemsep=3pt,topsep=2pt]
\item A real qualification profile should test transient recovery, trip exposure, compute feasibility, and standard-case RMSE.
\end{itemize}
\vspace{2pt}{\small\color{ofbBlue}$\rightarrow$ \itshape Reported in a standard, comparable form.}\\[2pt]
\vspace{3pt}
\begin{tcolorbox}[colback=ofbBlue!8,colframe=ofbBlue,boxrule=0.6pt,arc=2pt,left=4pt,right=4pt,top=3pt,bottom=3pt]
{\small\bfseries Key line: }{\small Qualify on all dimensions together, not one at a time.}
\end{tcolorbox}
\end{minipage}
\clearpage
\noindent\begin{minipage}[t]{0.585\textwidth}\vspace{0pt}
{\setlength{\fboxsep}{0pt}\setlength{\fboxrule}{0.4pt}\fbox{\includegraphics[width=\linewidth]{build/guion/g-57.png}}}
\end{minipage}\hfill
\begin{minipage}[t]{0.395\textwidth}\vspace{2pt}
{\footnotesize\color{ofbMuted}Slide 57 / 60 \quad $\sim$20s \quad cum 24:30}\hfill{\footnotesize\bfseries\color{ofbBlue}FIG}\\[2pt]
{\large\bfseries\color{ofbBlue} Standardization profile}\\[6pt]
{\small\itshape\color{ofbMuted}Say:}\\[2pt]
\begin{itemize}[leftmargin=1.1em,itemsep=3pt,topsep=2pt]
\item We propose reporting these dimensions explicitly, so any two estimators can be compared on equal terms.
\end{itemize}
\vspace{2pt}{\small\color{ofbBlue}$\rightarrow$ \itshape And when you truly need one number, here it is.}\\[2pt]
\vspace{3pt}
\begin{tcolorbox}[colback=ofbBlue!8,colframe=ofbBlue,boxrule=0.6pt,arc=2pt,left=4pt,right=4pt,top=3pt,bottom=3pt]
{\small\bfseries Key line: }{\small Make the comparison reproducible and fair.}
\end{tcolorbox}
\end{minipage}
\clearpage
\noindent\begin{minipage}[t]{0.585\textwidth}\vspace{0pt}
{\setlength{\fboxsep}{0pt}\setlength{\fboxrule}{0.4pt}\fbox{\includegraphics[width=\linewidth]{build/guion/g-58.png}}}
\end{minipage}\hfill
\begin{minipage}[t]{0.395\textwidth}\vspace{2pt}
{\footnotesize\color{ofbMuted}Slide 58 / 60 \quad $\sim$22s \quad cum 24:52}\hfill{\footnotesize\bfseries\color{ofbBlue}FIG}\\[2pt]
{\large\bfseries\color{ofbBlue} IBR Event Qualification Score}\\[6pt]
{\small\itshape\color{ofbMuted}Say:}\\[2pt]
\begin{itemize}[leftmargin=1.1em,itemsep=3pt,topsep=2pt]
\item When a single number is unavoidable, we define a composite, weighted, event-class score.
\end{itemize}
\vspace{2pt}{\small\color{ofbBlue}$\rightarrow$ \itshape Finally, where this work is going.}\\[2pt]
\vspace{3pt}
\begin{tcolorbox}[colback=ofbBlue!8,colframe=ofbBlue,boxrule=0.6pt,arc=2pt,left=4pt,right=4pt,top=3pt,bottom=3pt]
{\small\bfseries Key line: }{\small One number when you need it, without hiding the dimensions.}
\end{tcolorbox}
\end{minipage}
\clearpage
\noindent\begin{minipage}[t]{0.585\textwidth}\vspace{0pt}
{\setlength{\fboxsep}{0pt}\setlength{\fboxrule}{0.4pt}\fbox{\includegraphics[width=\linewidth]{build/guion/g-59.png}}}
\end{minipage}\hfill
\begin{minipage}[t]{0.395\textwidth}\vspace{2pt}
{\footnotesize\color{ofbMuted}Slide 59 / 60 \quad $\sim$22s \quad cum 25:14}\hfill{\footnotesize\bfseries\color{ofbTeal}CORE}\\[2pt]
{\large\bfseries\color{ofbTeal} Toward realtime IBR estimators}\\[6pt]
{\small\itshape\color{ofbMuted}Say:}\\[2pt]
\begin{itemize}[leftmargin=1.1em,itemsep=3pt,topsep=2pt]
\item Next: real-time estimators, integrating Simulink microgrid fault records, adding modern and machine-learning methods, C++ acceleration, and co-simulation with ANDES.
\end{itemize}
\vspace{2pt}{\small\color{ofbBlue}$\rightarrow$ \itshape Let me close.}\\[2pt]
\vspace{3pt}
\begin{tcolorbox}[colback=ofbTeal!8,colframe=ofbTeal,boxrule=0.6pt,arc=2pt,left=4pt,right=4pt,top=3pt,bottom=3pt]
{\small\bfseries Key line: }{\small This is where the benchmark goes next.}
\end{tcolorbox}
\end{minipage}
\clearpage
\noindent\begin{minipage}[t]{0.585\textwidth}\vspace{0pt}
{\setlength{\fboxsep}{0pt}\setlength{\fboxrule}{0.4pt}\fbox{\includegraphics[width=\linewidth]{build/guion/g-60.png}}}
\end{minipage}\hfill
\begin{minipage}[t]{0.395\textwidth}\vspace{2pt}
{\footnotesize\color{ofbMuted}Slide 60 / 60 \quad $\sim$35s \quad cum 25:49}\hfill{\footnotesize\bfseries\color{ofbTeal}CORE}\\[2pt]
{\large\bfseries\color{ofbTeal} Validation agenda and conclusion}\\[6pt]
{\small\itshape\color{ofbMuted}Say:}\\[2pt]
\begin{itemize}[leftmargin=1.1em,itemsep=3pt,topsep=2pt]
\item To conclude: there is no universal frequency estimator.
\item The defensible choice depends on the event, the metric, and the compute budget.
\item And OpenFreqBench gives the community a reproducible way to make that choice.
\end{itemize}
\vspace{3pt}
\begin{tcolorbox}[colback=ofbTeal!8,colframe=ofbTeal,boxrule=0.6pt,arc=2pt,left=4pt,right=4pt,top=3pt,bottom=3pt]
{\small\bfseries Key line: }{\small Thank you very much. I would be glad to take your questions.}
\end{tcolorbox}
\end{minipage}
\clearpage
\end{document}